\scanpage{158}
on l'appellera la \uline{section de Manin de} $F$ associée à l'op. de Picard-Fuchs $D$ et aux données $(P',\sigma)$.

Équation de Pic. Fuchs universelle : à valeur dans le pro-faisceau $P^\infty_{S/T}(E'')/\unsure{u|v}\delta u(E')$ ; elle redonne la section $\varphi$ du début.
% Two pairs of horizontal rules below this note.

Lorsque $E$ est seulem\textsuperscript{t} muni d'une connexion intégrable \deleted{\unsure{, de}} \inserted{rel. à $T$} (au lieu d'une stratification), on voit que les connexions intégrables sur $P=P'\mathbin{\overset{E'}{\times}}E$ correspondent aux hom
\[
 P'\xrightarrow{d_{P'}}\deleted{\unsure{E'}}E\otimes\Omega_{S/T} \tag*{(43)}
\]
satisfaisant
\[
 \left\{\begin{aligned}
 d_{P'}(x'+{\unsure{\xi|\zeta}}')&=d_{P'}(x')+d^0_E(u(\xi'))\\
 d_{P'}(P')&\subset Z^1(E\otimes\Omega_{S/T}),\ \text{i.e.}\quad d^{\unsure{1|2}}_E\circ d_{P'}=0 .
 \end{aligned}\right. \tag*{(44)}
\]
Dans ce cas, on a une variante de la notion d'équation \inserted{de} Picard-Fuchs \inserted{pour (30)}, en considérant les opérateurs différentiels
\[
 E\otimes\Omega^1_{S/T}\xrightarrow{\Delta\deleted{\unsure{D}}}F \tag*{(45)}
\]
nuls sur \deleted{l'image de $E'\otimes\Omega_{S/T}$} \inserted{$d^0_E(u(E'))$}, \deleted{i.e. les opérateurs différentiels sur} [qu'il suffira de prendre modulo ceux de la forme $\Delta'\circ d^1_E$, \deleted{$E''\otimes\Omega_{S/T}\longrightarrow$} où \deleted{$\Delta'$} $E\otimes\Omega^2_{S/T}\mathrel{\overset{\Delta'}{\text{---}}}F$ un autre op. diff]. Comme $d_{P'}$ (43) fournit \inserted{loc\textsuperscript{t}} une section de $E\otimes\Omega_{S/T}$ modulo une section
% archive sheet number: 157
